\documentclass{article}

% Language setting
% Replace `english' with e.g. `spanish' to change the document language
\usepackage[english]{babel}

% Set page size and margins
% Replace `letterpaper' with `a4paper' for UK/EU standard size
\usepackage[letterpaper,top=2cm,bottom=2cm,left=3cm,right=3cm,marginparwidth=1.75cm]{geometry}

% Useful packages
\usepackage{amsmath}
\usepackage{graphicx}
\usepackage[colorlinks=true, allcolors=blue]{hyperref}

\title{Labwork 1: Gradient Descent}
\author{Pham Quang Duong - 2540048}

\begin{document}
\maketitle


\section{Implement Gradient Descent}
\begin{itemize}
    \item[-] \textbf{Step 1:} Define f(x) and f\_(x)
    \item[-] \textbf{Step 2:} Radom initialize $x = x_0$, r, and thresh\_hold to stop
    \item[-] \textbf{Step 3:} Define the gradient\_descent() function with: 
    \begin{itemize}
        \item[$\bullet$]Assign x = x - r * f\_(x)
        \item[$\bullet$]Recompute f(x). Stop if f(x) is small enough or repeat if it's not
    \end{itemize}
\end{itemize}

\section{Effect of different learning rate r}
\begin{itemize}
    \item \textbf{Small learning rate:} The algorithm takes very small steps. This means that it requires many iterations to find the minimum value.
    
    \item \textbf{Good learning rate:} The algorithm works efficiently. It moves towards the minimum point quickly and reaches the correct answer in a reasonable number of steps.
    
    \item \textbf{Large learning rate:} The algorithm takes too many steps. It may jump over the minimum point between two values. In this case, it cannot find the minimum.
    
    \item \textbf{Very large learning rate:} The algorithm fails because it moves further away from the minimum in every step. 
\end{itemize}
\end{document}
